\section{The Dirac Equation}

\subsection{The Lorentz Group (Quick Facts)}
Lorentz transformations preserve the Minkowski metric:
\begin{equation}
\eta_{\rho\sigma}\Lambda^\rho_{\ \mu}\Lambda^\sigma_{\ \nu}=\eta_{\mu\nu}.
\end{equation}
Infinitesimally,
\begin{equation}
\Lambda^\mu_{\ \nu}=\delta^\mu_{\ \nu}+\omega^\mu_{\ \nu}, 
\qquad \omega_{\mu\nu}=-\omega_{\nu\mu}.
\end{equation}
Generators in the vector representation satisfy the Lorentz algebra:
\begin{equation}
[J^{\mu\nu},J^{\rho\sigma}]
=
i\left(
\eta^{\nu\rho}J^{\mu\sigma}
-\eta^{\mu\rho}J^{\nu\sigma}
-\eta^{\nu\sigma}J^{\mu\rho}
+\eta^{\mu\sigma}J^{\nu\rho}
\right).
\end{equation}

\subsection{Clifford Algebra and Gamma Matrices}
The Dirac gamma matrices $\gamma^\mu$ satisfy the Clifford algebra:
\begin{equation}
\{\gamma^\mu,\gamma^\nu\}=2\eta^{\mu\nu}\mathbf{1}.
\end{equation}
Define
\begin{equation}
\gamma^5 := i\gamma^0\gamma^1\gamma^2\gamma^3,
\end{equation}
which obeys
\begin{equation}
(\gamma^5)^2=\mathbf{1}, 
\qquad
\{\gamma^5,\gamma^\mu\}=0.
\end{equation}

Useful matrices:
\begin{equation}
\sigma^{\mu\nu}:=\frac{i}{2}[\gamma^\mu,\gamma^\nu].
\end{equation}

\subsection{Spinor Representation of Lorentz Transformations}
Dirac spinors transform as
\begin{equation}
\psi(x)\to \psi'(x') = S(\Lambda)\psi(x),
\qquad x'=\Lambda x,
\end{equation}
with
\begin{equation}
S(\Lambda)=\exp\left(-\frac{i}{4}\omega_{\mu\nu}\sigma^{\mu\nu}\right).
\end{equation}
Consistency condition:
\begin{equation}
S^{-1}(\Lambda)\gamma^\mu S(\Lambda)=\Lambda^\mu_{\ \nu}\gamma^\nu.
\end{equation}

\subsection{The Dirac Equation and Dirac Lagrangian}
The Dirac equation:
\begin{equation}
(i\gamma^\mu\partial_\mu-m)\psi=0.
\end{equation}
Dirac Lagrangian:
\begin{equation}
\mathcal{L}_D=\bar{\psi}(i\gamma^\mu\partial_\mu-m)\psi,
\qquad
\bar{\psi}:=\psi^\dagger\gamma^0.
\end{equation}
Euler--Lagrange gives
\begin{equation}
(i\gamma^\mu\partial_\mu-m)\psi=0,
\qquad
\bar{\psi}(i\overleftarrow{\partial}_\mu\gamma^\mu+m)=0.
\end{equation}

\paragraph{Klein--Gordon relation.}
Multiply the Dirac operator:
\begin{equation}
(i\gamma^\mu\partial_\mu-m)(i\gamma^\nu\partial_\nu+m)\psi
=
(-\partial^2-m^2)\psi=0,
\end{equation}
so each component satisfies the KG equation.

\subsection{Chiral Projectors and Chiral Spinors}
Define projectors:
\begin{equation}
P_L=\frac{1-\gamma^5}{2},\qquad P_R=\frac{1+\gamma^5}{2}.
\end{equation}
Decompose:
\begin{equation}
\psi=\psi_L+\psi_R,
\qquad
\psi_{L,R}=P_{L,R}\psi.
\end{equation}
For $m\neq 0$, left and right components mix in the Dirac equation.
For $m=0$, they decouple.

\subsection{The Weyl Equation (Massless Spinors)}
For $m=0$:
\begin{equation}
i\gamma^\mu\partial_\mu\psi=0.
\end{equation}
Chiral components satisfy separate Weyl equations:
\begin{equation}
i\gamma^\mu\partial_\mu\psi_L=0,
\qquad
i\gamma^\mu\partial_\mu\psi_R=0.
\end{equation}
Interpretation: massless fermions can have definite chirality.

\subsection{Parity}
Parity acts on coordinates:
\begin{equation}
\mathsf{P}:\quad (t,\mathbf{x})\to (t,-\mathbf{x}).
\end{equation}
On spinors (standard convention):
\begin{equation}
\psi(t,\mathbf{x})\xrightarrow{\mathsf{P}}
\psi'(t,\mathbf{x})=\gamma^0\psi(t,-\mathbf{x}).
\end{equation}
Under parity, chirality flips:
\begin{equation}
\psi_L \leftrightarrow \psi_R.
\end{equation}
Thus parity is incompatible with purely chiral theories (e.g. weak interaction).

\subsection{Majorana Spinors (Idea)}
A Majorana spinor satisfies a reality condition:
\begin{equation}
\psi^c=\psi,
\end{equation}
where
\begin{equation}
\psi^c := C\bar{\psi}^{\,T},
\end{equation}
and $C$ is the charge conjugation matrix (representation-dependent).

\paragraph{Physical meaning.}
A Majorana fermion is its own antiparticle.
Majorana mass terms schematically look like
\begin{equation}
\mathcal{L}_M \sim -\frac{m}{2}\left(\bar{\psi}\psi^c+\bar{\psi}^c\psi\right),
\end{equation}
and require neutral fermions (no conserved $U(1)$ charge).

\subsection{Symmetries and Conserved Currents}
\paragraph{Global $U(1)$ phase symmetry.}
\begin{equation}
\psi \to e^{i\alpha}\psi,
\qquad
\bar{\psi}\to \bar{\psi}e^{-i\alpha}.
\end{equation}
Noether current:
\begin{equation}
j^\mu=\bar{\psi}\gamma^\mu\psi,
\qquad
\partial_\mu j^\mu=0.
\end{equation}
Conserved charge:
\begin{equation}
Q=\int d^3x\, j^0.
\end{equation}

\paragraph{Axial (chiral) current.}
For $m=0$, the Lagrangian is also invariant under
\begin{equation}
\psi \to e^{i\alpha\gamma^5}\psi,
\end{equation}
with axial current
\begin{equation}
j_5^\mu=\bar{\psi}\gamma^\mu\gamma^5\psi.
\end{equation}
Classically,
\begin{equation}
\partial_\mu j_5^\mu = 0 \quad (m=0),
\end{equation}
but quantum mechanically this symmetry can be broken by anomalies (important in gauge theories).

\subsection{Plane Wave Solutions}
Seek solutions of the form
\begin{equation}
\psi(x)=u(p)e^{-ip\cdot x},
\qquad p^0=E_{\mathbf{p}}.
\end{equation}
Plugging into Dirac equation gives
\begin{equation}
(\slashed{p}-m)u(p)=0,
\qquad
\slashed{p}:=\gamma^\mu p_\mu.
\end{equation}
Similarly for negative-frequency solutions:
\begin{equation}
\psi(x)=v(p)e^{ip\cdot x},
\qquad
(\slashed{p}+m)v(p)=0.
\end{equation}

\paragraph{Spin sums (very useful).}
\begin{equation}
\sum_s u^{(s)}(p)\bar{u}^{(s)}(p)=\slashed{p}+m,
\qquad
\sum_s v^{(s)}(p)\bar{v}^{(s)}(p)=\slashed{p}-m.
\end{equation}

\paragraph{Normalization (common covariant choice).}
\begin{equation}
\bar{u}^{(s)}(p)u^{(r)}(p)=2m\,\delta^{sr},
\qquad
\bar{v}^{(s)}(p)v^{(r)}(p)=-2m\,\delta^{sr}.
\end{equation}

\subsection{Summary (Exam Checklist)}
\begin{itemize}
\item Gamma matrices satisfy $\{\gamma^\mu,\gamma^\nu\}=2\eta^{\mu\nu}$.
\item Dirac equation: $(i\slashed{\partial}-m)\psi=0$.
\item Lorentz transformations on spinors: $S(\Lambda)=\exp(-\frac{i}{4}\omega_{\mu\nu}\sigma^{\mu\nu})$.
\item Chirality projectors: $P_{L,R}=(1\mp\gamma^5)/2$.
\item Weyl spinors: massless chiral components decouple.
\item Parity: $\psi(t,\mathbf{x})\to \gamma^0\psi(t,-\mathbf{x})$, flips chirality.
\item Conserved current: $j^\mu=\bar{\psi}\gamma^\mu\psi$ from global $U(1)$.
\item Plane waves: $(\slashed{p}\mp m)u,v=0$ and spin sums $\sum u\bar{u}=\slashed{p}+m$.
\end{itemize}
